% SPDX-License-Identifier: Apache-2.0
% 03_from_netlist.tex: Drawing the netlist that synthesis writes
\section{Drawing the Netlist}
\label{sec:from-netlist}

A hand-drawn diagram shows the intent. The netlist shows what the tools
actually built. The two can differ in ways that matter: a register that
synthesis removed, a bus that got wider, or a carry chain where an adder was
expected. A drawing of the netlist shows those differences directly.

\subsection{From HDL to a JSON netlist}

The counter is synthesized by rules\_openxc7 \cite{rulesopenxc7}, which runs
the open flow for AMD 7-series parts. Yosys \cite{yosys} maps the design to
the part's primitives with \code{synth\_xilinx}. nextpnr \cite{openxc7} and
Project X-Ray \cite{prjxray} then place and route it, but the drawing needs
only the synthesis step. Its output is a JSON netlist: every cell, every
port, and the nets between them, numbered bit by bit.

\subsection{From the netlist to a drawing}

The command \code{blockpic import} reads the netlist and writes a
\code{.blockpic} file. From there, the netlist is drawn in exactly the same
way as Fig.~\ref{fig:architecture}. Table~\ref{tab:mapping} lists how each
part of the netlist is drawn.

\begin{table}[t]
\centering
\caption{How \code{blockpic import} draws a netlist.}
\label{tab:mapping}
\footnotesize
\begin{tabular}{@{}>{\raggedright}p{0.38\columnwidth}
                >{\raggedright\arraybackslash}p{0.54\columnwidth}@{}}
\toprule
In the netlist & In the drawing \\
\midrule
A cell, such as \code{FDRE} & A block of that kind \\
A module port & An input or output port block \\
A bus that goes whole to one pin & One thick wire \\
A bus whose bits go apart & A pin per bit, such as \code{O\_1} \\
A constant 0 or 1 & A \code{GND} or \code{VCC} block \\
Any other constant & A \code{CONST} block with its value \\
\bottomrule
\end{tabular}
\end{table}

\begin{figure*}[t]
\centering
\includegraphics[width=0.72\textwidth]{netlist.pdf}
\caption{The counter after synthesis for the xc7a200t, drawn by blockpic
from the netlist that Yosys wrote.}
\label{fig:netlist}
\end{figure*}

\subsection{Reading the drawing}

Fig.~\ref{fig:netlist} is the counter as Yosys built it. Each part of
Listing~\ref{lst:counter} can be found in it:
\begin{itemize}
\item The register is four \code{FDRE} flip-flops, one per bit. Each is
  drawn as a box with its pin names inside, as UG474 draws one. Their
  \code{R} pins are the synchronous reset, and \code{CE} is tied to
  \code{VCC}.
\item The adder became one \code{CARRY4}, the carry chain of a 7-series
  slice \cite{ug474}. Its \code{O\_1} to \code{O\_3} outputs feed the upper
  three flip-flops. Bit 0 needs no adder: an \code{INV} toggles it.
\item Every port went through a buffer: \code{IBUF} on \code{clk} and
  \code{reset}, \code{OBUF} on each LED. The clock also goes through a
  \code{BUFG} onto the global clock network.
\item The bus \code{led} was split into \code{led\_0} to \code{led\_3},
  because each bit leaves through its own buffer.
\end{itemize}

None of this is in the hand-drawn Fig.~\ref{fig:architecture}, and none of it
should be. The two figures answer different questions. The architecture
says what the counter does. The netlist says what the FPGA will run.

\begin{pitfall}
A netlist drawing grows with the design. A counter fits a column, but a
processor gives thousands of cells and an unreadable picture. Synthesis for
the FPGA also flattens the design into one module. To draw one part of a
large design, synthesize that part on its own, with it as the top level.
\end{pitfall}
